% =============================================================================
%  preamble.tex — Incompletitud, formalizada
%  Estilo, paquetes y los cinco entornos del libro.
%  Ver PLAN-LIBRO.md §4. NO contiene texto del libro.
% =============================================================================

% --- Motor -------------------------------------------------------------------
% El libro se compila con LuaLaTeX. Motivo: el proyecto usa Unicode
% intensivamente (σ ⟹ ⊢ ≐ ⌜·⌝ ⟨·⟩ Σ₁ Δ₀ ˙) y fontspec + luaotfload permiten
% una cadena de fuentes de respaldo por glifo. Con pdflatex esto no es viable.
\usepackage{iftex}
\ifLuaTeX\else\ifXeTeX\else
  \PackageError{preamble}{Compila con LuaLaTeX o XeLaTeX, no con pdfLaTeX}{%
    El libro necesita fontspec para el Unicode del proyecto.}%
\fi\fi

\usepackage[a4paper,margin=2.6cm,bottom=3cm]{geometry}
\usepackage{fontspec}
% babel español si la instalación de TeX lo trae; si no, se compila igual y
% sólo se pierden la partición de palabras y los rótulos automáticos.
\IfFileExists{spanish.ldf}{%
  \usepackage[spanish,es-noshorthands]{babel}%
}{%
  \PackageWarningNoLine{preamble}{spanish.ldf no disponible:
    instala texlive-lang-spanish para partición de palabras correcta}%
  \AtBeginDocument{%
    \renewcommand{\contentsname}{Índice}%
    \renewcommand{\partname}{Parte}%
    \renewcommand{\chaptername}{Capítulo}%
    \renewcommand{\appendixname}{Apéndice}%
    \renewcommand{\bibname}{Bibliografía}%
    \renewcommand{\tablename}{Tabla}%
    \renewcommand{\figurename}{Figura}%
  }%
}
\usepackage{csquotes}
\usepackage{microtype}
\usepackage{xcolor}
\usepackage{amsmath,amssymb,amsthm,mathtools}
\usepackage{booktabs}
\usepackage{tabularx}
\usepackage{array}
\usepackage{enumitem}
\usepackage{etoolbox}
\usepackage[most]{tcolorbox}
\usepackage{fancyvrb}
\usepackage{fvextra}  % breaklines para fancyvrb
\usepackage{graphicx}
\usepackage{hyperref}
\usepackage[nameinlink]{cleveref}

% Margen de estiramiento para párrafos con mucho \texttt: sin él, una racha de
% identificadores sin puntos de corte desborda por unos pocos puntos.
\setlength{\emergencystretch}{2.5em}

% --- Fuentes -----------------------------------------------------------------
% Texto: se prefiere Libertinus (buena cobertura griega y matemática);
% si no está instalada, Latin Modern.
\IfFontExistsTF{Libertinus Serif}{%
  \setmainfont{Libertinus Serif}\setsansfont{Libertinus Sans}%
}{\IfFontExistsTF{Latin Modern Roman}{%
  \setmainfont{Latin Modern Roman}\setsansfont{Latin Modern Sans}%
}{%
  \setmainfont{DejaVu Serif}\setsansfont{DejaVu Sans}%
}}

% Código: DejaVu Sans Mono cubre casi todo el Unicode del proyecto.
% Se declara una CADENA DE RESPALDO por glifo: si DejaVu no tiene el signo,
% luaotfload lo busca en las siguientes. Sin esto aparecen cuadraditos en
% ⌜·⌝, ≐ o los puntos de «dotar» (ȧ).
% Sólo LuaLaTeX admite cadena de respaldo por glifo (luaotfload). Bajo XeLaTeX
% se depende de la cobertura de la fuente elegida.
\makeatletter
\newif\ifrpp@fb  % ¿hay cadena de respaldo por glifo?
\ifLuaTeX
\IfFileExists{luaotfload-main.lua}{%
  \directlua{
    if luaotfload and luaotfload.add_fallback then
      luaotfload.add_fallback("leanfallback", {
        "DejaVuSansMono:mode=harf;",
        "Symbola:mode=harf;",
        "NotoSansMath-Regular:mode=harf;",
        "FreeMono:mode=harf;",
      })
    end
  }%
  \rpp@fbtrue
}{}
\fi
\IfFontExistsTF{DejaVu Sans Mono}{%
  \ifrpp@fb
    \setmonofont{DejaVu Sans Mono}[Scale=MatchLowercase,
                                   RawFeature={fallback=leanfallback}]%
  \else
    \setmonofont{DejaVu Sans Mono}[Scale=MatchLowercase]%
  \fi
}{%
  \setmonofont{Latin Modern Mono}[Scale=MatchLowercase]%
}
\makeatother

% --- Paleta ------------------------------------------------------------------
\definecolor{rppTinta}   {HTML}{1B1B1F}  % texto
\definecolor{rppCodigo}  {HTML}{F6F7F9}  % fondo de código
\definecolor{rppBorde}   {HTML}{C9CED6}
\definecolor{rppMate}    {HTML}{2F5D8C}  % matematica  — azul sobrio
\definecolor{rppLeccion} {HTML}{1F7A63}  % leccion     — verde azulado
\definecolor{rppMuro}    {HTML}{B0641E}  % muro        — ámbar: obstrucción
\definecolor{rppRefut}   {HTML}{9B2226}  % refutado    — rojo: resultado negativo
\definecolor{rppSondeo}  {HTML}{6E6E78}  % marca «fuera del build»
\definecolor{rppKw}      {HTML}{7A3E9D}
\definecolor{rppCom}     {HTML}{6E7781}
\definecolor{rppStr}     {HTML}{1F7A63}

\hypersetup{colorlinks=true, linkcolor=rppMate, citecolor=rppMate,
            urlcolor=rppMate, linktoc=all}

% --- Resaltado de Lean 4 -----------------------------------------------------
% NO se usa `listings`. Medido el 2026-09-03: con Unicode, listings REORDENA los
% caracteres —`(φ : Formula)` sale impreso como `φ( : Formula)`—, es decir,
% ALTERA EN SILENCIO el código publicado. En un libro cuyo principio editorial
% §2.1 es justo que eso no ocurra, es descalificante. Tampoco `minted`: exige
% --shell-escape y un lexer de Lean 4 en Pygments que no se puede dar por dado.
%
% Se usa `fancyvrb`, que es fiel byte a byte, y el RESALTADO LO HACE EL
% EXTRACTOR: scripts/extraer.py emite \kw{}/\cm{}/\st{} ya marcados. Es además
% lo correcto arquitectónicamente: quien conoce el código es el extractor.
\newcommand{\kw}[1]{\textcolor{rppKw}{\textbf{#1}}}   % palabra clave
\newcommand{\cm}[1]{\textcolor{rppCom}{\textit{#1}}}  % comentario
\newcommand{\st}[1]{\textcolor{rppStr}{#1}}           % cadena

% =============================================================================
%  MARCAS DE NIVEL — objeto vs meta, siempre a la vista
% =============================================================================
% Por qué esto no es decoración. La inconsistencia de ADR-012 fue un ERROR DE
% CATEGORÍA: `tcFn` es una operación sobre SINTAXIS declarada como función
% OBJETO. Si el nivel no se ve, el error no se ve. De ahí que cada fragmento de
% código y cada enunciado del libro lleve su nivel impreso.
%
% Y no es una distinción binaria: hay DOS EJES independientes.
%   · ¿QUIÉN afirma?  Lean (metateoría) · ⊢ (cálculo ω) · Prf (cálculo finitario)
%   · ¿SOBRE QUÉ?     Lean · sintaxis · objeto · código · Prov
% `prf_formCode_numeral` afirma en Prf SOBRE códigos; `hFN` afirma en ⊢ sobre lo
% mismo; `two_mul_consN` afirma en Lean sobre Nat. Son tres casillas distintas.

\definecolor{nvLean}  {HTML}{5B6470}   % metateoría: Lean
\definecolor{nvDer}   {HTML}{2F5D8C}   % el cálculo ⊢ (ω)
\definecolor{nvPrf}   {HTML}{1F7A63}   % el cálculo Prf (finitario, clásico)
\definecolor{nvPrf0}  {HTML}{3E9C84}   % Prf₀ (finitario, intuicionista)
\definecolor{nvPrfH}  {HTML}{15594A}   % PrfH (finitario CON contexto)
\definecolor{nvObj}   {HTML}{B0641E}   % lenguaje objeto
\definecolor{nvCod}   {HTML}{7A3E9D}   % códigos (términos objeto que denotan sintaxis)
\definecolor{nvProv}  {HTML}{9B2226}   % dentro de Prov

\makeatletter
\newcommand{\rpp@nvcolor}[1]{%
  \ifcsname nvcol@#1\endcsname \csname nvcol@#1\endcsname \else nvLean\fi}
\expandafter\def\csname nvcol@Lean\endcsname{nvLean}
\expandafter\def\csname nvcol@derives\endcsname{nvDer}
\expandafter\def\csname nvcol@Prf\endcsname{nvPrf}
\expandafter\def\csname nvcol@Prf0\endcsname{nvPrf0}
\expandafter\def\csname nvcol@PrfH\endcsname{nvPrfH}
\expandafter\def\csname nvcol@sintaxis\endcsname{nvLean}
\expandafter\def\csname nvcol@objeto\endcsname{nvObj}
\expandafter\def\csname nvcol@codigo\endcsname{nvCod}
\expandafter\def\csname nvcol@Prov\endcsname{nvProv}

\newcommand{\rpp@nvlabel}[1]{%
  \ifcsname nvlab@#1\endcsname \csname nvlab@#1\endcsname \else #1\fi}
\expandafter\def\csname nvlab@Lean\endcsname{Lean}
% La etiqueta lleva el TIPO a la vista: si el juicio arrastra contexto (Γ) o no.
% Confundir un juicio con contexto y uno sin él es un error de tipo, no de estilo.
\expandafter\def\csname nvlab@derives\endcsname{$\Gamma\vdash$}
\expandafter\def\csname nvlab@Prf\endcsname{Prf}
\expandafter\def\csname nvlab@Prf0\endcsname{Prf$_0$}
\expandafter\def\csname nvlab@PrfH\endcsname{$\Gamma\,$PrfH}
\expandafter\def\csname nvlab@sintaxis\endcsname{sintaxis}
\expandafter\def\csname nvlab@objeto\endcsname{objeto}
\expandafter\def\csname nvlab@codigo\endcsname{código}
\expandafter\def\csname nvlab@Prov\endcsname{Prov}

% \nivel{quién afirma}{sobre qué} — la chapa de dos casillas.
\newcommand{\nivel}[2]{%
  \begingroup\sffamily\scriptsize\bfseries
  \colorbox{\rpp@nvcolor{#1}}{\textcolor{white}{\strut\,\rpp@nvlabel{#1}\,}}%
  \colorbox{\rpp@nvcolor{#2}!18}{\textcolor{\rpp@nvcolor{#2}}{\strut\,\rpp@nvlabel{#2}\,}}%
  \endgroup}
\makeatother

% Marcas de nivel en prosa y en fórmulas.
\newcommand{\obj}[1]{\textcolor{nvObj}{#1}}        % símbolo del lenguaje objeto
\newcommand{\cod}[1]{\textcolor{nvCod}{#1}}        % código
\newcommand{\prov}[1]{\textcolor{nvProv}{#1}}      % dentro de Prov

% --- Nomenclatura (PLAN-LIBRO.md §2.7) --------------------------------------
% \defterm marca LA introducción de un término. El control `make terminos`
% exige que ningún término del vocabulario aparezca en el texto ANTES de su
% \defterm — sin exenciones: un lector con bases tiene que poder avanzar.
\newcommand{\defterm}[1]{\textbf{\textcolor{rppMate}{#1}}}

% Etiqueta de nivel de un TÉRMINO del glosario. Contesta, para cada palabra, la
% pregunta que el libro no puede dejar implícita: ¿esto es objeto o es meta?
\makeatletter
\newcommand{\nivelterm}[1]{%
  \begingroup\sffamily\scriptsize\bfseries
  \ifcsname ntcol@#1\endcsname
    \colorbox{\csname ntcol@#1\endcsname!16}{\textcolor{\csname ntcol@#1\endcsname}{\strut\,#1\,}}%
  \else
    \colorbox{nvLean!16}{\textcolor{nvLean}{\strut\,#1\,}}%
  \fi
  \endgroup}
\expandafter\def\csname ntcol@meta\endcsname{nvLean}
\expandafter\def\csname ntcol@objeto\endcsname{nvObj}
\expandafter\def\csname ntcol@meta→objeto\endcsname{nvCod}
\expandafter\def\csname ntcol@ambos\endcsname{nvDer}
\expandafter\def\csname ntcol@proyecto\endcsname{nvPrf}
\makeatother
\newcommand{\nota}[1]{{\footnotesize\color{rppCom} #1}}

% =============================================================================
%  LOS CINCO ENTORNOS
%  Ver PLAN-LIBRO.md §4. Cada uno tiene un papel distinto y no intercambiable.
% =============================================================================

% --- (1) lean — código Lean 4 real ------------------------------------------
% Uso normal: NO se escribe a mano. El extractor genera extraido/<clave>.tex y
% el capítulo lo trae con \leanfrag{clave} (§2.1: nada de copiar-pegar).
% El argumento opcional es la referencia de origen (módulo:declaración).
\newtcolorbox{leanbox}[1][]{%
  enhanced, breakable, sharp corners,
  colback=rppCodigo, colframe=rppBorde, boxrule=0.4pt,
  left=2mm, right=2mm, top=1.4mm, bottom=1.4mm,
  fonttitle=\ttfamily\scriptsize, coltitle=rppTinta,
  attach boxed title to top left={xshift=2mm, yshift=-2mm},
  boxed title style={colback=rppBorde, sharp corners, boxrule=0pt},
  #1}

\DefineVerbatimEnvironment{lean}{Verbatim}{%
  fontsize=\small, commandchars=\\\{\},
  breaklines=true, breakanywhere=false, breakautoindent=true,
  % El símbolo de continuación NO puede ser una flecha: pdftotext la extrae como
  % `→`, que es un token legítimo de Lean — un lector podría leer una flecha de
  % más en una línea partida, y el verificador del PDF no podría distinguirlas.
  % Se usa «›», que no aparece ni una vez en el código del proyecto ni en FOL.
  breaksymbolleft=\textcolor{rppBorde}{\guilsinglright},
  breaksymbolright={},
  xleftmargin=0pt, samepage=false}

% Trae un fragmento ya extraído del repo.
\newcommand{\leanfrag}[1]{\input{extraido/#1.tex}}

% Referencia de origen impresa bajo un fragmento.
\newcommand{\leanorigen}[2]{%
  \par\vspace{-0.4\baselineskip}%
  {\ttfamily\scriptsize\color{rppCom} #1\ \textbf{·}\ #2\par}}

% Footprint de `#print axioms` — OBLIGATORIO en todo teorema (§2.2(c)).
\newcommand{\footprint}[1]{%
  \par\vspace{-0.2\baselineskip}%
  {\ttfamily\scriptsize\color{rppCom}\#print axioms $\rightarrow$ #1\par}}

% Marca «fuera del build» para fragmentos de sondeos/ (§2.3).
\newcommand{\fueradelbuild}[1]{%
  \par\vspace{-0.2\baselineskip}%
  {\sffamily\scriptsize\color{rppSondeo}%
   \fcolorbox{rppSondeo}{white}{\strut\,fuera del build\,}\quad
   recompilar con \texttt{lake env lean #1}\par}}

% --- (2) matematica — el enunciado formal, legible sin Lean ------------------
% Sin barra de título: el rótulo va como línea de texto, para que el enunciado
% se lea como un enunciado y no como un aviso.
\NewDocumentEnvironment{matematica}{O{}}{%
  \begin{tcolorbox}[enhanced, breakable, sharp corners,
    colback=white, colframe=white, boxrule=0pt,
    borderline west={1.6pt}{0pt}{rppMate},
    left=4mm, right=2mm, top=1.6mm, bottom=1.6mm]%
  \ifblank{#1}{}{{\sffamily\bfseries\small\color{rppMate} #1\par}\vspace{0.4em}}%
}{%
  \end{tcolorbox}%
}

% --- (3) leccion — la nota pedagógica ---------------------------------------
\newtcolorbox{leccion}[1][Lección]{%
  enhanced, breakable, sharp corners,
  colback=rppLeccion!4, colframe=rppLeccion, boxrule=0.5pt,
  left=3mm, right=3mm, top=1.8mm, bottom=1.8mm,
  fonttitle=\bfseries\sffamily\small, coltitle=white,
  colbacktitle=rppLeccion,
  attach boxed title to top left={xshift=3mm, yshift=-2.4mm},
  boxed title style={sharp corners, boxrule=0pt},
  title={#1}}

% --- (4) muro — una obstrucción encontrada ----------------------------------
\newtcolorbox{muro}[1][Muro]{%
  enhanced, breakable, sharp corners,
  colback=rppMuro!5, colframe=rppMuro, boxrule=0.5pt,
  left=3mm, right=3mm, top=1.8mm, bottom=1.8mm,
  fonttitle=\bfseries\sffamily\small, coltitle=white,
  colbacktitle=rppMuro,
  attach boxed title to top left={xshift=3mm, yshift=-2.4mm},
  boxed title style={sharp corners, boxrule=0pt},
  title={#1}}

% --- (5) refutado — un resultado NEGATIVO, compilado ------------------------
% No es un muro: un muro es «no sé pasar»; esto es «se ha probado que no se
% pasa». Es el entorno que sostiene el capítulo 25 y media Parte IV.
\newtcolorbox{refutado}[1][Refutado]{%
  enhanced, breakable, sharp corners,
  colback=rppRefut!5, colframe=rppRefut, boxrule=0.5pt,
  left=3mm, right=3mm, top=1.8mm, bottom=1.8mm,
  fonttitle=\bfseries\sffamily\small, coltitle=white,
  colbacktitle=rppRefut,
  attach boxed title to top left={xshift=3mm, yshift=-2.4mm},
  boxed title style={sharp corners, boxrule=0pt},
  title={#1}}

% =============================================================================
%  Entornos de teorema (notación matemática estándar)
% =============================================================================
\theoremstyle{plain}
\newtheorem{teorema}{Teorema}[chapter]
\newtheorem{lema}[teorema]{Lema}
\newtheorem{proposicion}[teorema]{Proposición}
\newtheorem{corolario}[teorema]{Corolario}
\theoremstyle{definition}
\newtheorem{definicion}[teorema]{Definición}
\newtheorem{convenio}[teorema]{Convenio}
\theoremstyle{remark}
\newtheorem*{observacion}{Observación}

% =============================================================================
%  Macros de prosa
% =============================================================================
% \detokenize: permite escribir \ident{ax_tc_cons} sin escapar el subrayado.
% (El nombre \lean está tomado por el ENTORNO `lean`; en prosa se usa \ident.)
\newcommand{\ident}[1]{\texttt{\detokenize{#1}}}         % identificador Lean
% Las rutas de módulo son largas y aparecen constantemente. Con \texttt son
% inquebrables y desbordan la caja de texto; \nolinkurl parte por «/» y «.»,
% que es donde una ruta se puede partir sin confundir al lector.
\newcommand{\modulo}[1]{\begingroup\urlstyle{tt}\nolinkurl{#1}\endgroup}
\newcommand{\adr}[1]{\textsc{adr}-#1}                    % referencia a un ADR
\newcommand{\printaxioms}{\texttt{\#print axioms}}       % \detokenize duplica #
\newcommand{\doc}[1]{\textsf{\detokenize{#1}}}           % fichero .md del repo
\newcommand{\codigo}[1]{\ensuremath{\ulcorner #1 \urcorner}}  % ⌜φ⌝
\newcommand{\dotado}[1]{\ensuremath{\dot{#1}}}           % ȧ  («dotar»)
\newcommand{\Prov}{\operatorname{Prov}}
\newcommand{\Prf}{\mathsf{Prf}}
\newcommand{\derives}{\vdash}
\newcommand{\noderives}{\nvdash}

% Aviso editorial recurrente (§6): se usa donde toque, no se improvisa.
\newcommand{\avisoreparacion}{%
  \begin{muro}[Aviso]
  Lo que se retiró fue la inconsistencia \textbf{conocida y localizada}
  (\ident{ax_tc_cons}). \textbf{Eso no es una prueba de consistencia} de
  Q\texttt{++} extendido.
  \end{muro}}
\newcommand{\avisomediagodel}{%
  \begin{muro}[Aviso]
  Gödel I está \textbf{a medias}: sólo $\noderives G$. La indecidibilidad
  ($\noderives \lnot G$) \textbf{no está cerrada} en la cadena real.
  \end{muro}}
